\section{Conclusion}
\label{sec:conclusion}

We set out to answer three questions that VERINA leaves open. On the same \NTASKS{} tasks and with the same model, the specification language turned out to matter most. Moving from Lean to Dafny to Python/iContract raises the share of programs that pass their tests from \R{etwo.verina.pass}\% to \R{etwo.dafny.pass}\% to \R{etwo.python.pass}\%, and Python stays well ahead of Dafny inside every agent loop, although part of that lead comes from CrossHair's weaker, bounded guarantee. Verifier feedback improves on a testing-only agent by \R{gap3.dafny.verifier_vs_testing.diff}~pp on Dafny. Its larger benefit is that the agent no longer stops early: the testing-only agent declared success on unverified code in \R{dafny.testing_only.falsesuccess}\% of its Dafny runs, and no verifier-based agent did. Combining the two signals works best. Tests show that a large share of verifier failures (\R{ambig.dafny.share}\% in our runs) are missing proofs rather than wrong code, and SpecLoop, which uses that information, is the most accurate configuration, although the margin over verifier-only feedback (\R{gap2.dafny.full_vs_verifier.diff}~pp) is smaller than we hypothesised.

The remaining failures point to clear next steps. Most failed Dafny runs never compile, so the agent needs better grounding in the language, for example retrieved idioms for the construct that failed. Most of the rest are proof gaps where the invariant is too weak, which is where dedicated invariant generators could be plugged into the router's \textsc{proof\_gap} branch. We would also like to repeat the study with stronger and open-weight models, which the budget did not allow here, and on the harder VERINA-ADV tasks. Because every model response is cached, adding a model or a task only costs the new calls.
